% ============================================================
%  Chương II – Nền tảng mật mã khóa công khai
%  ~600 từ · 2 trang
% ============================================================
\chapter{Cơ sở Mật mã học và Hệ thống Khóa}
\label{chap:crypto}

% Lời dẫn chương: Mở đầu bằng một bức tranh tổng thể
Nền móng của an toàn thông tin được xây dựng trên những viên gạch toán học phức tạp của Mật mã học (Cryptography).
Tuy nhiên, lịch sử phát triển của ngành mật mã không phải là một đường thẳng, mà là một quá trình tiến hóa liên tục để giải quyết những nghịch lý phát sinh từ chính các giải pháp trước đó. Sự tiến hóa này bắt nguồn từ mật mã đối xứng với tốc độ vượt trội nhưng vướng mắc ở khâu phân phối khóa; dẫn đến sự ra đời của mật mã bất đối xứng nhằm giải quyết bài toán trao đổi khóa nhưng lại tạo ra lỗ hổng về xác thực danh tính.
Hạ tầng Khóa công khai (PKI) xuất hiện như một mảnh ghép hoàn thiện, kết nối sức mạnh toán học với cơ chế quản lý niềm tin trong thế giới thực. Chương này sẽ trình bày chi tiết chuỗi logic nền tảng đó.

% ----------------------------------------------------------
\section{Mật mã đối xứng và hạn chế}
\label{sec:2.1}

\subsubsection*{Nguyên lý hoạt động của Mật mã đối xứng}

Mật mã đối xứng (Symmetric-key Cryptography), hay còn gọi là mật mã khóa bí mật, là hình thức mã hóa cổ điển và lâu đời nhất trong lịch sử nhân loại. Nguyên lý cốt lõi của hệ thống này nằm ở tính "đối xứng" của chìa khóa: người gửi và người nhận sử dụng \textbf{cùng một khóa duy nhất ($K$)} cho cả hai quá trình mã hóa (Encryption) và giải mã (Decryption).

Về mặt toán học, quá trình này có thể được biểu diễn thông qua hai hàm cơ bản:
\begin{align*}
    C &= \mathrm{E}(K, P) \\
    P &= \mathrm{D}(K, C)
\end{align*}
Trong đó, $P$ là bản rõ (Plaintext), $C$ là bản mã (Ciphertext), hàm $\mathrm{E}$ và $\mathrm{D}$ lần lượt là thuật toán mã hóa và giải mã sử dụng khóa bí mật $K$.

Ưu điểm tuyệt đối và lý do khiến mật mã đối xứng vẫn là trụ cột không thể thay thế trong truyền tải dữ liệu hiện đại chính là \textbf{tốc độ xử lý (High Performance)}. Các thuật toán đối xứng tiêu biểu như \textbf{DES} (Data Encryption Standard), \textbf{3DES} (Triple DES) và đặc biệt là \textbf{AES} (Advanced Encryption Standard) được thiết kế dựa trên các phép toán mức bit cực kỳ đơn giản đối với phần cứng máy tính (như phép XOR, dịch chuyển bit Shift, và hoán vị Permutation). Thay vì phải giải quyết các bài toán toán học phức tạp như hệ mật mã bất đối xứng, AES xử lý dữ liệu theo từng khối (Block cipher) với tốc độ có thể lên tới hàng Gigabytes trên giây, đáp ứng hoàn hảo cho việc truyền phát Video, mã hóa ổ cứng hay các luồng dữ liệu thời gian thực.

\subsubsection*{Bài toán phân phối khóa (The Key Distribution Problem)}

Mặc dù sở hữu hiệu năng xuất sắc, mật mã đối xứng gặp khó khăn cấu trúc trong triển khai trên quy mô lớn: \textbf{Bài toán phân phối khóa}.

Nghịch lý của hệ thống đối xứng nằm ở một vòng luẩn quẩn (Catch-22): Để hai bên Alice và Bob có thể giao tiếp an toàn qua mạng Internet, họ cần phải chia sẻ trước với nhau một chiếc khóa $K$. Tuy nhiên, để gửi chiếc khóa $K$ này cho nhau một cách an toàn mà không bị kẻ thứ ba nghe lén, họ lại cần một kênh liên lạc đã được bảo mật từ trước. Nếu họ đã có sẵn một kênh bảo mật, thì họ đâu cần thiết lập hệ thống mật mã làm gì nữa? Trong quá khứ, giải pháp duy nhất là trao đổi khóa qua một kênh ngoại tuyến (out-of-band channel) như gặp mặt trực tiếp hoặc gửi qua người đưa thư (courier), điều này là bất khả thi trong kỷ nguyên kết nối toàn cầu.

Bài toán phân phối khóa tạo ra một rào cản khủng khiếp về mặt mở rộng quy mô (Scalability). Trong một mạng lưới có $n$ thực thể (node), nếu mỗi cặp thực thể đều muốn liên lạc bảo mật với nhau, hệ thống sẽ yêu cầu số lượng khóa phải quản lý tăng theo cấp số nhân. Công thức tính tổng số cặp khóa riêng biệt trong mạng là:
$$ \text{Tổng số khóa} = \frac{n(n-1)}{2} = O(n^2) $$

Hãy thử làm một phép tính thực tế:
\begin{itemize}
    \item Nếu mạng nội bộ chỉ có 10 nhân viên, số khóa cần quản lý là $10 \times 9 / 2 = 45$ khóa. (Có thể quản lý thủ công).
    \item Nếu một hệ thống ngân hàng có 1,000 khách hàng, số khóa cần tạo ra là gần $500,000$ khóa.
    \item Nếu triển khai trên phạm vi Internet với 1 triệu người dùng, chúng ta sẽ cần tới xấp xỉ \textbf{500 tỷ} chiếc khóa bí mật khác nhau.
\end{itemize}
Việc lưu trữ, phân phối, cập nhật và thu hồi hàng tỷ chiếc khóa bí mật này tạo ra một cơn ác mộng về mặt quản trị mạng, khiến mô hình đối xứng thuần túy trở nên hoàn toàn vô dụng đối với môi trường mạng phi tập trung.

\subsubsection*{Giải pháp KDC (Key Distribution Center) và những giới hạn chí mạng}

Để khắc phục bài toán $O(n^2)$, mô hình \textbf{Trung tâm Phân phối Khóa (KDC)} được đề xuất. Thay vì chia sẻ khóa với mọi thực thể, mỗi người dùng chỉ cần thiết lập một khóa chủ duy nhất với KDC (đóng vai trò Bên thứ ba đáng tin cậy). KDC sẽ sinh ra và phân phối Khóa phiên (Session Key) cho từng phiên giao tiếp, giúp thu hẹp số lượng khóa hệ thống cần quản lý xuống còn $O(n)$.

Tuy nhiên, KDC mang trong mình hai \textbf{giới hạn chí mạng} khiến nó không thể trở thành chuẩn mực cho môi trường Internet mở:

\begin{itemize}
    \item \textbf{Điểm lỗi tập trung (SPOF) và Rủi ro bảo mật:} KDC là trung tâm điều phối mọi luồng liên lạc và nắm giữ toàn bộ khóa chủ. Nếu KDC bị sập, toàn bộ mạng lưới tê liệt. Nghiêm trọng hơn, nếu KDC bị xâm nhập, kẻ tấn công có khả năng nghe lén và giả mạo mọi giao tiếp trong hệ thống.
    \item \textbf{Vấn đề ranh giới niềm tin (Trust Boundary):} KDC yêu cầu sự tin tưởng tuyệt đối vào một cơ quan trung tâm. Mô hình này hoạt động hiệu quả trong mạng nội bộ (như Kerberos), nhưng hoàn toàn bế tắc trên Internet mở, nơi không tồn tại một bên thứ ba nào đủ uy tín để các cá nhân và tổ chức toàn cầu đồng loạt giao phó bí mật.
\end{itemize}

Sự bế tắc của KDC trong việc thiết lập niềm tin qua biên giới đòi hỏi một cơ chế cho phép hai bên xa lạ giao tiếp bảo mật mà không cần chia sẻ trước khóa bí mật chung. Đó chính là động lực để \textbf{Mật mã khóa bất đối xứng} bước ra ánh sáng.

% ----------------------------------------------------------
\section{Mật mã khóa công khai}
\label{sec:2.2}

\subsubsection*{Cặp khóa bất đối xứng và Trapdoor function}

Khác với hệ mật mã đối xứng chỉ sử dụng một khóa duy nhất, mật mã khóa công khai (Public-key Cryptography) hoạt động dựa trên một cặp khóa gắn kết chặt chẽ về mặt toán học: \textbf{Khóa công khai (Public Key - PU)} và \textbf{Khóa bí mật (Private Key - PR)}. Tùy thuộc vào cách phối hợp và thứ tự sử dụng hai chiếc khóa này, hệ thống có thể giải quyết ba bài toán cốt lõi của an toàn thông tin:

\begin{enumerate}
    \item \textbf{Đảm bảo tính bảo mật (Confidentiality):} Người gửi mã hóa thông điệp bằng $PU$ của người nhận. Lúc này, bản mã chỉ có thể được giải mã bằng $PR$ duy nhất mà người nhận đang cất giữ. Kịch bản này bảo vệ dữ liệu khỏi những kẻ nghe lén trên đường truyền.
    \item \textbf{Xác thực và Chữ ký số (Authentication/Digital Signature):} Người gửi "mã hóa" (hoặc ký) thông điệp bằng chính $PR$ của mình. Bất kỳ ai trên mạng cũng có thể dùng $PU$ của người gửi để giải mã. Nếu quá trình giải mã thành công, điều này minh chứng tuyệt đối rằng thông điệp thực sự xuất phát từ người gửi đó và nội dung không hề bị chỉnh sửa (chống chối cãi).
    \item \textbf{Bảo mật và Xác thực song hành:} Người gửi kết hợp cả hai thao tác. Đầu tiên, ký thông điệp bằng $PR$ của nguồn gửi để xác thực, sau đó tiếp tục bọc thêm một lớp mã hóa bằng $PU$ của đích đến để bảo mật. Nhờ vậy, thông điệp vừa được giữ kín, vừa đảm bảo tính chính danh.
\end{enumerate}

% --- KHU VỰC CHÈN 3 ẢNH SƠ ĐỒ ---
\begin{figure}[htbp]
    \centering
    \begin{subfigure}[b]{\textwidth}
        \centering
        \slidefig[0.88]{EWPU.png}
        \caption{Mã hóa bằng khóa công khai của người nhận}
    \end{subfigure}
    \par\vspace{0.5cm}
    \begin{subfigure}[b]{\textwidth}
        \centering
        \slidefig[0.88]{EWPR.png}
        \caption{Xác thực / chữ ký bằng khóa riêng người gửi}
    \end{subfigure}
    \par\vspace{0.5cm}
    \begin{subfigure}[b]{\textwidth}
        \centering
        \slidefig[0.88]{EWB.png}
        \caption{Bảo mật và xác thực kết hợp}
    \end{subfigure}
    \caption{Ba kịch bản ứng dụng cốt lõi của hệ mật mã khóa bất đối xứng}
    \label{fig:asymmetric_scenarios}
\end{figure}
% ---------------------------------------

\paragraph*{Kịch bản 1: Đảm bảo tính bảo mật (Encryption with Public Key)}
Mục tiêu của kịch bản này là đảm bảo chỉ duy nhất người nhận hợp pháp mới có thể đọc được nội dung thông điệp (Confidentiality).
\begin{itemize}
    \item \textbf{Phía người gửi (Bob):} Bob lấy thông điệp gốc $X$ và tra cứu Khóa công khai của người nhận (Alice) trên một vòng khóa (public key ring). Sau đó, Bob thực hiện mã hóa:
    $$Y = \mathrm{E}[PU_a, X]$$
    Bản mã $Y$ lúc này được truyền đi trên kênh không an toàn.
    \item \textbf{Phía người nhận (Alice):} Khi nhận được $Y$, Alice sử dụng Khóa bí mật duy nhất của mình ($PR_a$) để giải mã:
    $$X = \mathrm{D}[PR_a, Y]$$
    Vì chỉ có Alice mới nắm giữ $PR_a$, mọi kẻ nghe lén dù bắt được $Y$ và biết $PU_a$ cũng không thể đảo ngược quá trình để tìm ra $X$.
\end{itemize}

\paragraph*{Kịch bản 2: Xác thực nguồn gốc và Chữ ký số (Encryption with Private Key)}
Mục tiêu ở đây không phải là giấu kín dữ liệu, mà là chứng minh tuyệt đối nguồn gốc của thông điệp (Authentication).
\begin{itemize}
    \item \textbf{Phía người gửi (Bob):} Bob dùng chính Khóa bí mật của mình ($PR_b$) để "mã hóa" (bản chất là ký số) lên thông điệp $X$:
    $$Y = \mathrm{E}[PR_b, X]$$
    \item \textbf{Phía người nhận (Alice):} Alice nhận $Y$ và tra cứu Khóa công khai của Bob ($PU_b$) để khôi phục:
    $$X = \mathrm{D}[PU_b, Y]$$
    Bất kỳ ai trên mạng cũng có thể dùng $PU_b$ để đọc được thông điệp này. Tuy nhiên, nếu hàm giải mã $\mathrm{D}$ chạy thành công và trả về dữ liệu có nghĩa, điều đó minh chứng 100\% rằng bản mã $Y$ \textit{bắt buộc} phải được tạo ra từ $PR_b$ của Bob. Bob không thể chối cãi việc mình đã gửi thông điệp này.
\end{itemize}

\paragraph*{Kịch bản 3: Bảo mật và Xác thực song hành (Secrecy and Authentication)}
Để đáp ứng cả hai yêu cầu khắt khe nhất trong liên lạc, hệ thống sẽ bọc dữ liệu qua hai lớp vỏ mã hóa, kết hợp cả hai kịch bản trên.
\begin{itemize}
    \item \textbf{Phía người gửi (Source A):} Đầu tiên, A tạo chữ ký để xác thực bằng cách mã hóa với Khóa bí mật của mình: $Y = \mathrm{E}[PR_a, X]$. Sau đó, A bọc lớp bảo mật bên ngoài bằng cách mã hóa tiếp $Y$ với Khóa công khai của đích đến B:
    $$Z = \mathrm{E}[PU_b, Y]$$
    \item \textbf{Phía người nhận (Destination B):} Quá trình bóc tách diễn ra ngược lại. B dùng Khóa bí mật của mình để phá lớp vỏ bảo mật: $Y = \mathrm{D}[PR_b, Z]$. Cuối cùng, B dùng Khóa công khai của A để xác thực và lấy ra thông điệp gốc:
    $$X = \mathrm{D}[PU_a, Y]$$
\end{itemize}

Để các kịch bản trên có thể hoạt động an toàn trong thực tế, hệ thống đòi hỏi một cơ chế toán học cực kỳ khắt khe. Việc công khai $PU$ và thuật toán mã hóa cho cả thế giới biết đặt ra thách thức: Làm sao để ngăn chặn hacker tính ngược lại để tìm ra $PR$ hoặc bản rõ ban đầu? Câu trả lời nằm ở khái niệm \textbf{Hàm một chiều có cửa sập}.

Khái niệm này bao gồm hai tính chất then chốt:
\begin{itemize}
    \item \textbf{Hàm một chiều (One-way Function):} Là một hàm số $Y = f(X)$ thỏa mãn điều kiện: nếu biết $X$, việc tính xuôi ra $Y$ là rất dễ dàng (thời gian đa thức). Tuy nhiên, nếu chỉ biết $Y$, việc tính ngược $X = f^{-1}(Y)$ là \textit{bất khả thi về mặt tính toán (computationally infeasible)} đối với máy tính hiện hành.
    \item \textbf{Cửa sập (Trapdoor):} Hệ mật mã không thể chỉ dùng hàm một chiều đơn thuần (vì chính người nhận cũng sẽ không giải mã được). Do đó, hàm này được thiết kế thêm một "cửa sập". Nếu sở hữu thông tin bổ sung $k$ (chính là Khóa bí mật $PR$), người nhận có thể kích hoạt cửa sập và dễ dàng tính toán ngược $X = f_k^{-1}(Y)$.
\end{itemize}

Tóm lại, sự sinh tồn của mật mã bất đối xứng phụ thuộc vào việc xây dựng thành công hàm tính toán một chiều vững chắc, nơi Khóa công khai đóng vai trò là hàm $f$ để chốt khóa, còn Khóa bí mật chính là chìa khóa cửa sập $k$ để mở luồng tính ngược.

\subsubsection*{Thuật toán RSA}

Thuật toán RSA, được đặt theo tên của ba nhà phát minh Rivest, Shamir và Adleman, là một trong những hệ mật mã khóa công khai phổ biến và nền tảng nhất. Tính an toàn của RSA dựa trên độ khó của \textbf{bài toán phân tích thừa số nguyên tố lớn (Integer Factorization Problem)} \cite{Rivest1978}. Về mặt toán học, việc nhân hai số nguyên tố khổng lồ để tạo ra một hợp số thì diễn ra rất nhanh, nhưng việc làm ngược lại — phân tích hợp số đó ra hai thừa số ban đầu — là một bài toán không khả thi tính toán (intractable) đối với khả năng của máy tính hiện tại.

\paragraph*{Quy trình hoạt động của RSA}
Hệ thống RSA được vận hành qua 3 giai đoạn cốt lõi:

\begin{enumerate}
    \item \textbf{Tạo khóa (Key Generation):} (Thực hiện bởi người nhận - Alice)
    \begin{itemize}
        \item Chọn hai số nguyên tố lớn ngẫu nhiên $p$ và $q$ ($p \neq q$).
        \item Tính modulus $n = p \times q$.
        \item Tính giá trị hàm Euler totient: $\phi(n) = (p-1)(q-1)$.
        \item Chọn một số nguyên $e$ sao cho $1 < e < \phi(n)$ và $\gcd(\phi(n), e) = 1$ (tức là $e$ và $\phi(n)$ nguyên tố cùng nhau).
        \item Tính khóa giải mã $d$ là nghịch đảo của $e$ theo modulo $\phi(n)$, tức là:
        $$d \equiv e^{-1} \pmod{\phi(n)}$$
        \item \textbf{Khóa công khai (Public Key):} $PU = \{e, n\}$.
        \item \textbf{Khóa bí mật (Private Key):} $PR = \{d, n\}$. \\
        \textit{(Lưu ý: Các giá trị $p, q$ và $\phi(n)$ phải được tiêu hủy hoặc cất giữ tuyệt mật sau khi tạo khóa xong).}
    \end{itemize}

    \item \textbf{Mã hóa (Encryption):} (Thực hiện bởi người gửi - Bob)
    \begin{itemize}
        \item Bob muốn gửi thông điệp $M$ cho Alice (với điều kiện $M < n$).
        \item Bob sử dụng Khóa công khai $PU$ của Alice để tính bản mã $C$:
        $$C = M^e \bmod n$$
    \end{itemize}

    \item \textbf{Giải mã (Decryption):} (Thực hiện bởi người nhận - Alice)
    \begin{itemize}
        \item Alice nhận bản mã $C$ và sử dụng Khóa bí mật $PR$ của mình để khôi phục nguyên trạng thông điệp $M$:
        $$M = C^d \bmod n$$
    \end{itemize}
\end{enumerate}

% --- KHU VỰC CHÈN ẢNH SƠ ĐỒ RSA ---
\begin{figure}[htbp]
    \centering
    \slidefigtall{RSA.png}
    \caption{Sơ đồ tạo khóa, mã hóa và giải mã của thuật toán RSA (Nguồn: Figure 9.5 \& 9.6)}
    \label{fig:rsa_algorithm}
\end{figure}
% ---------------------------------------

\paragraph*{Tính bảo mật của RSA}

Độ an toàn của hệ mật mã RSA phải đối mặt với 4 hướng tấn công chủ đạo từ giới phân tích mật mã (cryptanalysis):
\begin{enumerate}
    \item \textbf{Tấn công vét cạn (Brute force):} Thử tất cả các khóa bí mật (Private Key) có thể. Hướng này bị vô hiệu hóa bằng cách sử dụng không gian khóa khổng lồ.
    \item \textbf{Tấn công toán học (Mathematical attacks):} Tập trung vào việc giải bài toán phân tích thừa số nguyên tố lớn để dịch ngược các thông số cốt lõi.
    \item \textbf{Tấn công thời gian (Timing attacks):} Một dạng tấn công kênh kề (side-channel), đo lường thời gian chạy của thuật toán giải mã để suy luận ra từng bit của khóa bí mật.
    \item \textbf{Tấn công bản mã lựa chọn (Chosen Ciphertext attacks - CCA):} Khai thác các tính chất toán học cơ bản của thuật toán RSA để tạo ra các thông điệp có chủ đích.
\end{enumerate}

\textbf{1. Bài toán phân tích thừa số (The Factoring Problem)} \\
Phòng tuyến toán học đầu tiên và quan trọng nhất của RSA là khó khăn trong việc phân tích một hợp số lớn $n$ thành hai thừa số nguyên tố $p$ và $q$. Nếu kẻ tấn công có thể phân tích thành công $n$, chúng sẽ dễ dàng tính được $\phi(n) = (p-1)(q-1)$ và từ đó trích xuất được khóa giải mã $d$. Mặc dù các thuật toán phân tích (như General Number Field Sieve - GNFS) ngày càng tiến bộ, nhưng nỗ lực tính toán vẫn tăng theo cấp số mũ so với độ dài của khóa. Do đó, việc nâng kích thước khóa RSA lên tiêu chuẩn 2048-bit hoặc 4096-bit được coi là an toàn tuyệt đối trước các đòn tấn công toán học truyền thống.

\textbf{2. Tấn công thời gian (Timing Attack) và Cách khắc phục} \\
Timing Attack là một kịch bản tấn công thực tế và cực kỳ nguy hiểm do Paul Kocher đề xuất. Kẻ tấn công không cần giải toán mà sẽ đo lường thời gian phản hồi (chậm/nhanh) của bộ xử lý khi thực thực thi thuật toán lũy thừa modulo (Square and Multiply). Sự chênh lệch thời gian xử lý giữa bit `0` và bit `1` của khóa $d$ sẽ tố cáo toàn bộ giá trị của khóa.

Để chống lại đòn tấn công này, có 3 biện pháp kỹ thuật được áp dụng:
\begin{itemize}
    \item \textbf{Cố định thời gian tính toán (Constant exponentiation time):} Buộc hệ thống luôn phải tiêu tốn một khoảng thời gian cố định cho mọi phép tính, bất kể bit đó là `0` hay `1`. Tuy nhiên, điều này làm giảm hiệu năng hệ thống.
    \item \textbf{Thêm độ trễ ngẫu nhiên (Random delay):} Chèn thêm các khoảng thời gian chờ ngẫu nhiên vào thuật toán để làm nhiễu thông số đo lường của kẻ tấn công.
    \item \textbf{Tung hỏa mù (Blinding):} Đây là biện pháp tối ưu nhất. Trước khi giải mã, bản mã $C$ được nhân với một giá trị ngẫu nhiên $r$: $C' = C \cdot r^e \bmod n$. Bộ xử lý sẽ giải mã trên một cục dữ liệu rác $C'$, khiến mọi phép đo thời gian của hacker trở nên vô nghĩa, sau đó hệ thống sẽ khử $r$ ở bước cuối cùng để trả về bản rõ $M$.
\end{itemize}

\textbf{3. Tấn công CCA và chuẩn OAEP (Optimal Asymmetric Encryption Padding)} \\
Thuật toán RSA nguyên thủy (Basic RSA) có một lỗ hổng nghiêm trọng dựa trên đặc tính nhân (multiplicative property): Mã hóa của $M_1$ nhân với mã hóa của $M_2$ sẽ bằng mã hóa của $(M_1 \times M_2)$. Kẻ tấn công có thể lợi dụng điều này bằng cách chèn thêm các hệ số nhân vào bản mã $C$, biến đổi nó thành một bản mã có chủ đích (Chosen Ciphertext) và lừa máy chủ giải mã để lấy thông tin.

Để vô hiệu hóa hoàn toàn các phương pháp tấn công này, hệ thống RSA thực tế không bao giờ mã hóa trực tiếp thông điệp $M$, mà phải đi qua một cơ chế "đệm" cực kỳ phức tạp có tên là \textbf{OAEP (Optimal Asymmetric Encryption Padding)}.

Về bản chất, OAEP sẽ trộn lẫn thông điệp $M$ ban đầu với các chuỗi padding, các giá trị băm $H(P)$ và một hạt giống ngẫu nhiên (Seed). Thông qua mạng lưới kết hợp các hàm sinh mặt nạ (Mask Generating Function - MGF) và phép toán XOR chéo, OAEP sẽ nhào nặn thông điệp thành một khối dữ liệu (Encoded Message - EM) hoàn toàn hỗn độn và không còn bất kỳ mối liên hệ toán học tuyến tính nào. Nếu kẻ tấn công cố tình nhân thêm hệ số vào bản mã, khi giải mã, cấu trúc OAEP sẽ bị phá vỡ, máy chủ sẽ phát hiện lỗi và từ chối trả về kết quả.

% --- KHU VỰC CHÈN ẢNH SƠ ĐỒ OAEP ---
\begin{figure}[htbp]
    \centering
    \slidefigtall{OAEP.png}
    \caption{Sơ đồ khối quá trình nhào nặn dữ liệu của chuẩn OAEP (Nguồn: Figure 9.10)}
    \label{fig:rsa_oaep}
\end{figure}
% ---------------------------------------
% ---------------------------------------

\subsubsection*{Trao đổi khóa Diffie-Hellman}
% Ý chính 1: Sứ mệnh và định nghĩa cơ bản
Diffie-Hellman (DH) là một trong những hệ thống mật mã khóa công khai (Public-key Cryptography) đầu tiên được công bố. Khác với tính đa dụng của RSA, thuật toán DH được thiết kế với một mục đích duy nhất và chuyên biệt: \textbf{Trao đổi khóa (Key Exchange)}. Sứ mệnh cốt lõi của DH là cho phép hai thực thể hoàn toàn xa lạ có thể đàm phán và thiết lập một Khóa bí mật chung (Shared Secret / Session Key) một cách an toàn thông qua một kênh truyền tải không bảo mật, nơi mọi thông điệp đều có thể bị kẻ tấn công nghe lén. Khóa chung này sau đó sẽ được sử dụng làm đầu vào cho các thuật toán mã hóa đối xứng (như AES) nhằm truyền tải dữ liệu ở tốc độ cao.

% Ý chính 2: Nền tảng toán học bảo vệ thuật toán
Về mặt toán học, độ an toàn của thuật toán Diffie-Hellman phụ thuộc hoàn toàn vào độ khó của \textbf{Bài toán Logarit rời rạc (Discrete Logarithm Problem)}. Thuật toán sử dụng số nguyên tố $p$ rất lớn và một căn nguyên thủy (primitive root) $a$ của nó. Trong không gian số học modulo $p$, việc tính toán xuôi $a^i \bmod p$ là một thao tác cực kỳ nhanh chóng. Tuy nhiên, thao tác tính toán ngược — tức là tìm lại số mũ bí mật $i$ khi chỉ biết kết quả $b$, cơ số $a$ và số modulo $p$ ($i = \log_a b \pmod p$) — là một bài toán bất khả thi về mặt thời gian (infeasible) đối với các hệ thống máy tính hiện tại.

\paragraph*{Các bước thực hiện thuật toán Diffie-Hellman}

Giả sử hai người dùng A (Alice) và B (Bob) muốn thiết lập một khóa bí mật chung qua một kênh truyền không an toàn. Thuật toán được thực hiện qua 3 bước cốt lõi:

\begin{enumerate}
    \item \textbf{Khởi tạo thông số công khai (Global Public Elements):} \\
    Trước khi trao đổi, Alice và Bob phải thống nhất hai thông số nền tảng (các giá trị này có thể truyền công khai):
    \begin{itemize}
        \item $q$: Một số nguyên tố rất lớn.
        \item $\alpha$: Một căn nguyên thủy (primitive root) của $q$, với điều kiện $\alpha < q$.
    \end{itemize}

    \item \textbf{Tạo khóa cá nhân và khóa công khai (Key Generation):}
    \begin{itemize}
        \item \textbf{Phía Alice:}
        \begin{itemize}
            \item Chọn một số nguyên ngẫu nhiên bí mật $X_A < q$ (đóng vai trò là Private Key).
            \item Tính giá trị trộn công khai $Y_A = \alpha^{X_A} \bmod q$ và gửi $Y_A$ cho Bob.
        \end{itemize}
        \item \textbf{Phía Bob:}
        \begin{itemize}
            \item Tương tự, chọn một số nguyên ngẫu nhiên bí mật $X_B < q$.
            \item Tính giá trị trộn công khai $Y_B = \alpha^{X_B} \bmod q$ và gửi $Y_B$ cho Alice.
        \end{itemize}
    \end{itemize}

    \item \textbf{Tính toán khóa chung bí mật (Calculation of Secret Key):} \\
    Sau khi nhận được giá trị $Y$ từ đối phương, hai bên tiến hành tính toán khóa phiên (Session Key) $K$ độc lập tại hệ thống của mình:
    \begin{itemize}
        \item Alice tính: $K = (Y_B)^{X_A} \bmod q$
        \item Bob tính: $K = (Y_A)^{X_B} \bmod q$
    \end{itemize}
\end{enumerate}
% --- ĐOẠN CODE CHÈN ẢNH SƠ ĐỒ VÀO ĐÂY ---
\begin{figure}[htbp]
    \centering
    \slidefig[0.92]{DH_protocol.png}
    \caption{Sơ đồ luồng trao đổi khóa Diffie-Hellman (Nguồn: Figure 10.2)}
    \label{fig:dh_exchange}
\end{figure}
% -----------------------------------------

\paragraph*{Phân tích độ an toàn trước kẻ tấn công}
Giả sử có một kẻ nghe lén (Eavesdropper) kiểm soát toàn bộ kênh truyền, kẻ này sẽ thu thập được đầy đủ 4 thông số công khai: $q, \alpha, Y_A$ và $Y_B$. Tuy nhiên, để tìm ra được khóa chung $K$, kẻ tấn công bắt buộc phải trích xuất được số mũ bí mật $X_A$ hoặc $X_B$. Quá trình này đòi hỏi kẻ tấn công phải giải phương trình:
$$X_B = \mathrm{dlog}_{\alpha, q}(Y_B)$$
Trong đó $\mathrm{dlog}$ là phép tính logarit rời rạc. Nếu tham số $q$ được chọn đủ lớn (theo tiêu chuẩn hiện tại là từ 2048-bit trở lên), sức mạnh điện toán hiện tại của nhân loại là không thể giải quyết bài toán tính ngược này trong thời gian khả thi (computationally infeasible).

\subsubsection*{Sơ đồ chữ ký số (Digital Signatures)}

Trong các hệ thống xác thực sử dụng khóa đối xứng (Symmetric Authentication), việc trao đổi thông điệp giữa hai bên Alice và Bob luôn tiềm ẩn những nguy cơ tranh chấp pháp lý mà các kỹ thuật mã hóa thông thường không thể giải quyết triệt để. Để hiểu rõ tại sao chúng ta cần đến Chữ ký số, hãy xem xét hai kịch bản rủi ro sau đây:
\textbf{1. Kịch bản giả mạo nội dung (Forgery by Receiver):} \\
Giả sử Alice và Bob chia sẻ một khóa bí mật chung để xác thực các giao dịch chuyển tiền. Bob (người nhận) có thể tự ý chỉnh sửa con số trong thông điệp từ 100 USD thành 1000 USD, sau đó dùng khóa chung đó để tạo lại mã xác thực mới. Khi có tranh chấp, Bob có thể khẳng định trước tòa rằng Alice đã gửi 1000 USD, và vì mã xác thực là hợp lệ, Alice không có bằng chứng nào để minh oan.

\textbf{2. Kịch bản chối bỏ trách nhiệm (Denial of Sending):} \\
Alice gửi một yêu cầu đặt lệnh mua cổ phiếu cho Bob. Tuy nhiên, sau đó thị trường lao dốc khiến Alice thua lỗ. Lúc này, Alice hoàn toàn có thể tuyên bố rằng mình chưa từng gửi thông điệp đó. Vì Bob cũng nắm giữ khóa bí mật chung, Alice có thể đổ lỗi rằng Bob đã tự tạo ra thông điệp để hãm hại mình.

% --- KHU VỰC CHÈN ẢNH SƠ ĐỒ CÁC TÌNH HUỐNG TRANH CHẤP ---
\begin{figure}[htbp]
    \centering
    \slidefigtall{digitroll.png}
    \caption{Mô hình tổng quát của quá trình ký số và các kịch bản tranh chấp phát sinh (Nguồn: Figure 13.1)}
    \label{fig:digital_signature_model}
\end{figure}
% ---------------------------------------

\paragraph*{Yêu cầu đối với một Chữ ký số tiêu chuẩn}
Từ những rủi ro trên, chúng ta rút ra kết luận rằng trong các tình huống mà sự tin tưởng giữa người gửi và người nhận không phải là tuyệt đối, chúng ta cần một cơ chế mạnh mẽ hơn xác thực thông thường. Một sơ đồ chữ ký số (Digital Signature) đúng nghĩa phải thỏa mãn các đặc tính sau:

\begin{itemize}
    \item \textbf{Xác thực tác giả và thời điểm:} Phải có khả năng kiểm chứng chính xác ai là người ký và chữ ký đó được tạo ra vào thời gian nào.
    \item \textbf{Xác thực nội dung:} Tại thời điểm ký, nội dung thông điệp phải được "khóa" lại. Mọi sự thay đổi sau đó đều phải làm cho chữ ký trở nên vô hiệu.
    \item \textbf{Khả năng kiểm chứng bởi bên thứ ba:} Đây là đặc tính quan trọng nhất để giải quyết tranh chấp. Một bên thứ ba độc lập (ví dụ: Trọng tài hoặc Tòa án) phải có khả năng kiểm tra tính hợp lệ của chữ ký mà không cần nắm giữ khóa bí mật của người ký.
\end{itemize}

Như vậy, về mặt chức năng, chữ ký số bao hàm cả chức năng xác thực thông thường nhưng được nâng cấp thêm khả năng \textbf{Chống chối cãi (Non-repudiation)}, biến nó thành một công cụ pháp lý quan trọng trong giao dịch điện tử.

\paragraph*{Yêu cầu kỹ thuật đối với một hệ thống Chữ ký số an toàn}
Dựa trên việc phân tích các hình thức tấn công và nguy cơ giả mạo (Attacks and Forgeries), chúng ta có thể thiết lập một bộ quy tắc khắt khe để đảm bảo tính sinh tồn của một sơ đồ chữ ký số trong thực tế. Các yêu cầu này không chỉ nhắm vào tính bảo mật mà còn chú trọng đến hiệu năng vận hành:

\begin{itemize}
    \item \textbf{Sự phụ thuộc vào nội dung (Message Dependency):} Chữ ký số phải là một chuỗi bit được tạo ra dựa trên chính nội dung của thông điệp được ký. Nếu thông điệp thay đổi dù chỉ 1 bit, chữ ký phải thay đổi hoàn toàn. Điều này ngăn chặn đòn tấn công \textit{Existential Forgery} (giả mạo tồn tại).
    \item \textbf{Tính duy nhất của người gửi (Sender Uniqueness):} Chữ ký phải sử dụng thông tin duy nhất và bí mật của người gửi (Private Key). Yêu cầu này là "lá chắn" cốt lõi để chống lại việc giả mạo (forgery) từ phía người nhận và ngăn chặn việc chối bỏ trách nhiệm (denial) từ phía người gửi.
    \item \textbf{Tính khả thi trong tạo lập (Ease of Production):} Việc tạo ra chữ ký số phải tương đối dễ dàng và nhanh chóng đối với người gửi hợp pháp, không gây quá tải cho tài nguyên hệ thống.
    \item \textbf{Tính dễ dàng trong xác thực (Ease of Verification):} Mọi bên thứ ba hoặc người nhận hợp pháp đều phải có khả năng kiểm tra và xác định tính đúng đắn của chữ ký một cách thuận tiện thông qua khóa công khai.
    \item \textbf{Kháng giả mạo về mặt tính toán (Computational Infeasibility):} Đây là yêu cầu quan trọng nhất để đối phó với \textit{Universal} và \textit{Selective Forgery}. Hệ thống phải đảm bảo rằng kẻ tấn công không thể:
    \begin{itemize}
        \item Tạo ra một thông điệp mới cho một chữ ký số đã tồn tại.
        \item Tạo ra một chữ ký giả mạo cho một thông điệp đã cho sẵn.
    \end{itemize}
    \item \textbf{Khả năng lưu trữ (Storage Practicality):} Bản sao của chữ ký số phải có kích thước tối ưu để có thể lưu trữ lâu dài nhằm phục vụ mục đích đối soát hoặc giải quyết tranh chấp pháp lý sau này.
\end{itemize}
Để thỏa mãn đồng thời các yêu cầu khắt khe này, các hệ thống hiện đại thường sử dụng một \textbf{hàm băm an toàn (Secure Hash Function)} kết hợp với các thuật toán mã hóa bất đối xứng. Hàm băm đảm bảo tính phụ thuộc nội dung và hiệu năng tính toán, trong khi mật mã bất đối xứng đảm bảo tính duy nhất và khả năng xác thực bởi bên thứ ba. Tuy nhiên, việc thiết kế chi tiết các giao thức này (như cách chèn thêm padding) cần được thực hiện cực kỳ cẩn trọng để tránh các lỗ hổng tinh vi như tấn công thời gian hay tấn công bản mã lựa chọn.

% ----------------------------------------------------------
\section{Tại sao mật mã bất đối xứng chưa đủ?}
\label{sec:2.3}

Mật mã khóa công khai (Asymmetric Cryptography) đã giải quyết được một trong những bài toán nan giải nhất của ngành mật mã học: Phân phối khóa bí mật qua một kênh truyền không an toàn mà không cần hai bên phải gặp mặt trực tiếp. Bằng các đặc tính toán học ưu việt như hàm một chiều có cửa sập (Trapdoor one-way function), hệ thống này cung cấp cả hai tính năng cốt lõi là bảo mật (Confidentiality) và chữ ký số (Digital Signature).

Tuy nhiên, sức mạnh toán học đơn thuần là chưa đủ để xây dựng một hệ thống an toàn trên quy mô toàn cầu. Lỗ hổng chết người của mật mã bất đối xứng không nằm ở các phương trình toán học, mà nằm ở lỗ hổng về "Niềm tin" (Trust) trong quá trình phân phối và định danh khóa. Khi các khóa công khai thuần túy (Raw Public Keys) được truyền tải trên một mạng lưới mở như Internet, chúng trở thành miếng mồi ngon cho các đòn tấn công có chủ đích.

\subsubsection*{Tấn công Man-in-the-Middle (MITM) trên Raw Public Key}

Để minh họa cho điểm yếu này, hãy xem xét kịch bản trao đổi khóa công khai nguyên thủy (Raw Public Key Exchange). Giả sử Alice muốn gửi một tài liệu mật cho Bob. Theo lý thuyết, Alice chỉ cần xin khóa công khai của Bob ($K_{\mathrm{pub\_Bob}}$) để mã hóa tài liệu. Tuy nhiên, trên một môi trường mạng bị kiểm soát bởi kẻ tấn công (Mallory), lý thuyết này hoàn toàn sụp đổ.

Quá trình tấn công Kẻ đứng giữa (Man-in-the-Middle) diễn ra âm thầm và chia cắt hoàn toàn kênh liên lạc của hai nạn nhân mà họ không hề hay biết. Mallory đóng vai trò là một trạm trung chuyển độc hại, thao túng toàn bộ luồng giao tiếp. Diễn biến chi tiết được trình bày trong Bảng~\ref{tab:mitm}.

\begin{table}[htbp]
    \centering
    \caption{Kịch bản tấn công MITM trên quá trình trao đổi khóa công khai thuần túy}
    \label{tab:mitm}
    \renewcommand{\arraystretch}{1.5}
    \begin{tabular}{|p{0.25\linewidth}|p{0.4\linewidth}|p{0.3\linewidth}|}
        \hline
        \rowcolor{gray!20}
        \textbf{Alice} & \textbf{Mallory (Kẻ tấn công)} & \textbf{Bob} \\ \hline
        1. Gửi yêu cầu: \textit{"Bob, gửi cho tôi $K_{\mathrm{pub\_Bob}}$"} & $\xrightarrow{\text{Đánh chặn yêu cầu}}$ \newline Chuyển tiếp yêu cầu đến Bob. & Nhận yêu cầu từ "Alice". \\ \hline
        & $\xleftarrow{\text{Đánh chặn } K_{\mathrm{pub\_Bob}}}$ \newline Lưu giữ $K_{\mathrm{pub\_Bob}}$. & 2. Phản hồi: Gửi $K_{\mathrm{pub\_Bob}}$ cho Alice. \\ \hline
        3. Nhận $K_{\mathrm{pub\_Mallory}}$. Đinh ninh đây là khóa của Bob. & $\xrightarrow{\text{Gửi } K_{\mathrm{pub\_Mallory}}}$ \newline Đóng giả Bob, gửi khóa công khai của chính mình cho Alice. & \\ \hline
        4. Mã hóa tài liệu mật $M$: \newline $C = \mathrm{Enc}(K_{\mathrm{pub\_Mallory}}, M)$ \newline Gửi $C$ cho Bob. & $\xrightarrow{\text{Đánh chặn bản mã } C}$ \newline Giải mã bằng khóa bí mật của Mallory: $M = \mathrm{Dec}(K_{\mathrm{priv\_Mallory}}, C)$. \newline \textit{(Mallory đã đọc/sửa được dữ liệu)} & \\ \hline
        & $\xrightarrow{\text{Mã hóa lại và Gửi}}$ \newline $C' = \mathrm{Enc}(K_{\mathrm{pub\_Bob}}, M)$ \newline Gửi $C'$ cho Bob. & 5. Nhận $C'$. \newline Giải mã bằng $K_{\mathrm{priv\_Bob}}$. \newline Đinh ninh dữ liệu đến từ Alice và hoàn toàn bảo mật. \\ \hline
    \end{tabular}
\end{table}

Hậu quả của đòn tấn công này là thảm họa. Mặc dù thuật toán mã hóa (ví dụ RSA) không hề bị bẻ khóa, nhưng tính bảo mật đã bị vô hiệu hóa hoàn toàn do Alice đã sử dụng "nhầm" chìa khóa để khóa cửa. Kịch bản này không chỉ đúng với việc mã hóa dữ liệu mà còn áp dụng tương tự với việc xác minh chữ ký số: Mallory có thể thay thế khóa công khai, từ đó tự do giả mạo chữ ký mang tên Bob mà Alice vẫn xác minh thành công.

\subsubsection*{Vấn đề Ràng buộc Danh tính (The Identity Binding Problem)}

Cội nguồn của đòn tấn công MITM nêu trên xuất phát từ một thiếu sót nền tảng của mật mã khóa công khai: \textbf{Sự vắng mặt của tính ràng buộc danh tính (Identity Binding)}.

Về bản chất toán học, một khóa công khai (như $K_{\mathrm{pub}}$) chỉ là một tập hợp các con số vô tri vô giác (ví dụ: một cặp số $(e, n)$ trong RSA, hay một tọa độ điểm $(x, y)$ trên đường cong Elliptic). Bản thân các con số này không chứa đựng bất kỳ siêu dữ liệu (metadata) nào để chứng minh nó thuộc về thực thể nào. Không có bất kỳ nhãn dán nào trên khóa công khai ghi rằng \textit{"Khóa này là tài sản hợp pháp của Bob, nhân viên ngân hàng X"}.

Trong thế giới thực, để chứng minh một bức ảnh thuộc về một cái tên, chúng ta cần Chứng minh nhân dân hoặc Hộ chiếu --- một tài liệu do cơ quan nhà nước có thẩm quyền cấp phép, đóng dấu giáp lai để \textbf{ràng buộc} (bind) khuôn mặt với đặc điểm nhân thân.

Trên không gian mạng, nếu Alice nhận được một chuỗi bit dài 2048-bit từ một email có tên "Bob", làm sao cô ấy có thể tin tưởng 100\% đó không phải là Mallory giả mạo? Nếu không có một cơ chế kiểm chứng độc lập, mọi giao dịch dựa trên mật mã bất đối xứng đều được xây dựng trên một nền móng cát: niềm tin mù quáng vào kênh truyền tải.

\paragraph*{Cầu nối tới Hạ tầng Khóa Công khai (PKI)}
Để mật mã bất đối xứng có thể vận hành trong thế giới thực, nơi mà mọi nút mạng đều có thể là kẻ thù, chúng ta bắt buộc phải có một "Sở tư pháp" kỹ thuật số. Cần có một cơ chế, một định dạng tiêu chuẩn, và một Bên thứ ba đáng tin cậy (Trusted Third Party) đứng ra bảo lãnh cho tính chính danh của các khóa công khai.

Họ phải tạo ra một "tờ giấy thông hành" kỹ thuật số, trong đó niêm phong chặt chẽ \textbf{Khóa công khai (Public Key)} cùng với \textbf{Danh tính (Identity)} của chủ sở hữu bằng một Chữ ký số không thể làm giả. Cấu trúc đó chính là \textbf{Chứng chỉ số (Digital Certificate)}, và hệ sinh thái quản lý toàn bộ vòng đời của các chứng chỉ này chính là lý do \textbf{Hạ tầng Khóa Công khai (PKI - Public Key Infrastructure)} ra đời. Đây sẽ là trọng tâm phân tích của Chương~\ref{chap:pki-arch}.
